%======================================================================
\section{허락 층의 시빌 노출}
\label{sec:sybil}

\begin{yoji}
허락 하나에는 수수료만 들고 토큰이 묶이지 않아요. 균등 재시작으로 허락 층을 걷는 점수에서 닫힌 농장은 자기 인구 몫의 몇 배를 모아요. 수정 분석의 엔도스랭크 그래프에서는 주소 26개짜리 별 모양이면 목표를 코호트 점수의 상위 10\%로 올려요(등록 그래프에서는 13개). 결합한다고 이 빈틈이 닫히지는 않는다는 것도 논리로 보여요. 이 모형은 W1 라벨을 안 뒤에 계산했고 어느 규칙에도 들어가지 않아요.
\end{yoji}

앞의 결과는 정직한 자료에 관한 것이에요. 평판 층은 속이는 데 드는 비용으로도 판단해야 해요. 이 절은 Gy\"{o}ngyi와 Garcia-Molina의 농장 결과~\cite{gyongyi2005alliances}를, 균등 재시작을 쓰는 엔도스랭크의 허락 걷기에 맞춰 수치로 재요. 그래프는 수정 분석의 2026년 6월 30일 동결 그래프(점 2,513,982개, 화살표 3,808,949개)이고, 등록 분석의 그래프를 견줘 보여요. 닫힌 농장의 한계(식~\eqref{eq:farm})는 이 층을 균등 재시작으로 걷는 걷기라면 무엇에나, $\lambda=1$의 C-PR에도 성립해요. 하지만 아래의 별 모양 수는 코호트의 엔도스랭크 점수에 견줘 읽어요. 이 모형은 W1 라벨을 안 뒤에 계산했고 어느 규칙에도 들어가지 않아요.

\subsection{화살표 하나의 비용}

스펜더를 허락해도 아무것도 맡겨 두지 않아요. 토큰이 묶이지 않고, 허락하는 주소는 그 토큰을 하나도 갖지 않아도 돼요. 상한이 있는 무게에서는 몇백 달러짜리 허락이 무제한 허락과 거의 같은 무게를 가져요. 그리고 행을 정규화하면, 나가는 화살표가 하나뿐인 점은 금액이 얼마든 그 화살표가 행 전체를 가져요. 그래서 공격자의 두 주소 사이의 화살표는 아무것도 위험에 내놓지 않고, 수수료만 들어요.

수정 그래프는 동결 직전 12개월 동안 정한 허락만 담아요. 그래서 허락은 정한 지 1년이 지나면 엔도스랭크에서도, C-PR의 허락 층에서도 빠져요. 따라서 농장은 적어도 1년에 한 번은 허락을 다시 내야 해요. 닫힌 농장 안에서는 엔도스랭크의 나이 할인이 상관없어요. 농장 주소의 행은 모두 농장 안에 머무는 화살표들끼리 정규화되기 때문이에요. 이런 기간을 $q$번 거치는 동안 닫힌 농장의 비용은 다음과 같아요.
\begin{equation}
C=q\,\bigl(|E_{\mathcal S}|\,c_{\text{gas}}+c_{\text{rev}}\bigr).
\label{eq:cost}
\end{equation}
$c_{\text{gas}}$는 허락 하나의 수수료이고, $c_{\text{rev}}$는 기간마다 허락을 거두거나 들키지 않으려고 드는 비용이에요. $c_{\text{gas}}$는 재지 않았어요. 값이 얼마든 비용은 화살표 수와 기간 수에 따라 선형으로 늘어요. 그래프를 널리 거래되는 다섯 토큰으로 한정해도 이 비용은 오르지 않아요. 그 가운데 하나를 허락하는 것이 다른 토큰을 허락하는 것보다 비싸지 않기 때문이에요. 다섯 토큰으로 한정해서 없어지는 것은, 실제로 움직인 적 없는 주인의 이름으로 토큰 계약이 써 넣은 로그예요.

\subsection{닫힌 농장이 모으는 것}

농장 $\mathcal S$의 모든 점이 나가는 화살표를 가지되 그 화살표가 모두 농장 안을 가리키고, 정직한 주인은 아무도 농장 주소를 허락하지 않는다고 해요. 매달린 점들이 가진 정상 몫을 $r_D$라고 써요. 매달린 점은 모두 정직한 점이에요. 균등한 $\boldsymbol\pi$로 식~\eqref{eq:pr}을 $\mathcal S$에 걸쳐 더하면 다음이 나와요.
\begin{equation}
\sum_{i\in\mathcal S} r_i=\frac{|\mathcal S|}{|V|}\cdot\frac{1-d+d\,r_D}{1-d}\le\frac{|\mathcal S|}{(1-d)|V|+d|\mathcal S|}.
\label{eq:farm}
\end{equation}
부등식은 $r_D\le 1-\sum_{i\in\mathcal S}r_i$를 써서 얻어요. 둘째 인수는 농장이 자기 인구 몫의 몇 배를 차지하는지를 나타내요. 매달린 점이 있으면 그 몫이 다시 나눠져 농장에도 닿으므로, 이 배수는 늘 1보다 커요.
\begin{itemize}
\item 수정 그래프에서 $d=0.75$, $0.85$, $0.95$일 때 $r_D=0.365$, $0.394$, $0.419$예요. 그래서 닫힌 농장은 자기 인구 몫의 2.10배, 3.23배, 8.96배를 차지해요.
\item 주소 20개가 다음 주소를 하나씩 허락하는 고리를 그래프에 넣어 보면, $d=0.85$에서 자기 몫의 3.23배를 차지해요.
\item 등록 그래프에서는 배수가 2.01, 2.96, 6.62예요.
\end{itemize}

\begin{puri}[왜 인구 몫보다 더 많이 모이나요?]
걷는 사람이 매달린 점에 도착하면 그 몫은 모든 점에 고르게 다시 나눠져요. 농장 주소도 그 일부를 받아요. 그런데 닫힌 농장에서는 화살표가 모두 농장 안을 가리키므로, 받은 몫이 순간 이동(확률 $1-d$)으로만 빠져나가요. 그래서 몫이 농장 안에 쌓이고, 매달린 점의 몫 $r_D$가 클수록, 그리고 $d$가 1에 가까울수록 더 많이 쌓여요.
\end{puri}

\subsection{몫을 한곳에 모으기}

식~\eqref{eq:farm}는 농장 전체의 몫을 묶을 뿐이고, 공격자는 그 몫을 한 점에 모을 수 있어요. 별 모양에서는 농장 주소 $f$개가 모두 목표 $t$를 허락하고, $t$도 그 주소들을 하나하나 허락해요. 순간 이동과 매달린 행에서 각 점이 받는 몫을 $c=(1-d+d\,r_D)/|V|$로 두면, 균형 식에서 다음이 나와요.
\begin{equation}
r_t=\frac{c\,(1+d\,f)}{1-d^{2}}.
\label{eq:star}
\end{equation}
\begin{itemize}
\item 수정 그래프에서 $c=1.93\times10^{-7}$이에요. 들어오는 화살표가 없는 모든 점의 공통 점수예요.
\item $f=10$이면 주소 11개와 기간마다 허락 20개로, 목표는 별을 넣은 그래프에서 $6.60\times10^{-6}$을 받아요. 식~\eqref{eq:star}의 값이에요. 코호트 점수의 중앙값($5.82\times10^{-7}$)보다 위이고, 90번째 백분위($1.56\times10^{-5}$)의 5분의 2쯤이에요.
\item 농장 주소를 하나 더할 때마다 기간마다 허락 2개를 더 내고 $d\,c/(1-d^2)$를 더 얻어요.
\end{itemize}
식~\eqref{eq:star}을 $f$에 대해 풀면 각 순위에 드는 비용이 나와요.
\begin{itemize}
\item 중앙값에는 별이 아예 필요 없어요. 서로를 허락하는 주소 둘만으로 이미 $c/(1-d)$에 이르러요.
\item 90번째 백분위에는 농장 주소 26개가 들어요.
\item 99번째 백분위($1.35\times10^{-3}$)에는 약 2,292개가 들어요.
\item 코호트 계약 가운데 가장 높은 점수(0.0278)에는 약 47,073개가 들어요. 이 수는 그렇게 큰 농장이 $|V|$를 키워 $c$를 낮추는 효과를 무시한 값이에요.
\item 등록 그래프에서는 이 수가 13, 1,253, 약 1,080,271이에요.
\end{itemize}
그래서 농장은 기간마다 허락 52개로 코호트 점수의 상위 10분의 1에 들어가요(등록 그래프에서는 26개). 그 수수료는 재지 않았어요. 반면 맨 꼭대기에는 천 배가 넘게 큰 농장이 필요해요. 앞선 계약들이 정상 몫의 큰 부분을 쥐고 있기 때문이에요.

\begin{puri}[별 모양 농장의 셈을 따라가 보면]
별 모양에서 목표 $t$는 농장 주소 $f$개에게서 허락을 받고, 그 주소들에게 다시 허락을 줘요. 걷는 사람은 $t$와 농장 주소 사이를 오가므로, 순간 이동으로 들어온 몫이 $t$에 거듭 돌아와 쌓여요. 식~\eqref{eq:star}에서 $f$가 하나 늘면 $t$의 점수는 일정한 양만큼 늘어요. 그래서 수정 그래프에서는 상위 10\%에 몇십 개의 주소로 닿지만, 이미 큰 몫을 쥔 맨 위 계약을 넘으려면 주소가 몇만 개나 필요해요.
\end{puri}

\subsection{결합 걷기에서도 같은 논리}

$\lambda>0$인 C-PR에서는 식~\eqref{eq:alpha}의 다시 정규화 때문에, 나가는 화살표가 모두 허락인 점은 그 화살표를 무게 1로 따라가요. 그래서 서로만 허락하고 송금을 하지 않는 농장은, 허락 걷기에서와 똑같이 C-PR에서도 닫혀 있어요. 식~\eqref{eq:farm}는 C-PR의 매달린 몫으로 그대로 들어맞아요. 다만 저자들은 그 몫을 계산하지도, C-PR의 그래프에서 농장을 돌려 보지도 않았어요. 그래서 위의 배수를 C-PR에 대해 주장하지 않아요. 이 논리는 결합이 빈틈을 닫지 못한다는 것만 보여 줘요.

상한을 둔 C-PR의 층에서, 정직한 주인에게서 산 허락 하나가 몫을 크게 옮기는 것은 그 허락이 그 주인의 다른 허락에 견줘 클 때뿐이에요. 그리고 상한만큼의 허락이면 스펜더는 그 주인에게서 적어도 $U$달러를 꺼낼 수 있어요. 반면 엔도스랭크의 상한 있는 층에서는, 새로 낸 몇백 달러짜리 허락 하나가 중앙값 주인의 행에서 68.3\%를 차지해요. 대부분 주인의 마지막 허락이 오래됐거나 작기 때문이에요.

\subsection{처방, 시험하지 않음}

Cheng과 Friedman이 찾은 비대칭이 해결책을 가리켜요~\cite{cheng2005sybilproof}. 트러스트랭크처럼~\cite{gyongyi2004trustrank} 뛰기와 매달린 점의 몫을 검증한 씨앗 집합으로만 보내면, 씨앗 밖의 농장은 정직한 주인에게서 오는 화살표로만 몫을 얻어요. 그리고 그 화살표는 돈을 주고 사야 해요. 저자들은 씨앗으로 제한한 재시작을 평가하지 않았어요. S-PR은 그런 처방이 아니에요. S-PR의 씨앗은 농장이 부풀릴 수 있는 허락 점수이기 때문이에요. 휴먼 패스포트~\cite{humanpassportnddocs} 같은 자격 증명 관문은 다른 질문, 곧 한 주소가 서로 다른 한 사람의 것인지를 다뤄요. 이 질문은 계약에는 바로 맞는 짝이 없고, 그런 관문에는 나름의 약점이 있어요~\cite{bordeianu2026sybil}. 이 분석은 모형에 바탕을 두고, 관측한 그래프에서 풀이 프로그램으로 점검했어요. 체인 위에서 실제로 공격해 보지는 않았어요.

%======================================================================
\section{타당성을 위협하는 것}
\label{sec:threats}

\begin{yoji}
스펜더 라벨은 같은 종류의 화살표가 새로 생기는 것이라 지금의 수를 따라가는 점수에 유리해요. 수정 분석(달러 금액, 다섯 토큰, 12개월 점수 기간)은 등록 결과를 안 뒤에 정했고, 등록은 제3자의 시간 도장이 없으며, 트레이더는 70개뿐이에요. 증거는 한 롤업, 한 걸음 에고 네트워크, 고정된 전체 무게에서 나왔어요.
\end{yoji}

\subsection{구성 타당도}

\begin{itemize}
\item 스펜더 라벨은 층을 이루는 것과 같은 종류의 화살표가 새로 생기는 것이에요. 그래서 미래의 연결을 재고, 지금의 수를 따라가는 점수에 유리해요. 둘 다 신용, 안전, 손실의 결과가 아니에요.
\item 트레이더 라벨은 이 자료에서 부도에 가장 가까운 결과예요. 하지만 한 거래소에서의 강제 청산은 대출의 부도가 아니고, W1 트레이더 70개 가운데 9개만 청산을 겪었어요.
\item 허락 화살표는 기준 시점까지의 마지막 허락이지, 지금 유효한 허락이 아니에요(\ref{sec:model}절). 수정 분석에서는 12개월 점수 기간 전에 정한 허락은 아무리 오래 유효해도 세지 않아요. 동결 블록의 계약 상태에서 허락을 읽으면, 새 쌍으로 센 것 가운데 몇 개가 다시 낸 옛 관계였는지 가릴 수 있어요.
\item 층의 무게가 중요해요. 계획대로의 원래 금액 층에서는 무제한 허락과 소수 자릿수가 큰 토큰이 행을 차지했고, 수정 분석의 달러 금액과 상한이 그것을 막아요.
\item 두 분석은 규칙 A와 B에서 같은 판정에 이르렀어요. 하지만 둘은 무게, 토큰, 점수 기간이 한꺼번에 달라서, 어느 하나를 훑어본 것이 아니에요. 엔도스랭크와 AWP의 상한 있는 무게를 쓰는 층처럼 다른 결합은 이 원고에서 평가하지 않았어요.
\end{itemize}

\subsection{내적 타당도}

\begin{itemize}
\item 규칙 A는 이 연구가 로그를 하나도 추출하기 전에 정했어요. 규칙 B의 가설은 W0의 탐색적 대비에서, 그 허용 폭은 W0의 구간에서 나왔어요.
\item 등록은 저자들 자신의 저장소에 있는 커밋이지, 독립 등록처의 항목이 아니에요. 등록이 W1 추출보다 먼저였다는 것은 제3자의 시간 도장이 아니라 커밋 시각과 작업 시각에 기대요(\ref{sec:gate}절).
\item 달러 금액, 다섯 토큰, 12개월 점수 기간을 쓰는 수정 분석은 모든 등록 결과를 안 뒤에 정했어요. 그래서 그 선택이 한쪽 결과에 유리했을 수 있어요. 두 사실이 그 위험을 줄이지만 없애지는 못해요. 기울기와 상한은 W0 동결까지 12개월의 사건 분위수에서 나와서 라벨 기간의 사건이 들어가지 않아요. 그리고 수정 분석은 규칙 A와 B에서 등록 분석과 같은 판정에 이르렀어요.
\item 수정 분석을 먼저 보고하는 것도 결과를 본 뒤의 선택이에요. 등록 분석은 규칙을 정할 때 시험하려던 분석으로서, 결과의 모든 표와 그림에 남아 있어요.
\item 빠진 계약을 0점으로 두는 규칙은 처리 과정의 선택이에요. S2는 이 규칙이 등록 분석의 F5를 동점 구조를 통해 바꾸지만, 판정의 어느 부분도 바꾸지 않는다는 것을 보여 줘요.
\end{itemize}

\subsection{통계적 결론 타당도}

\begin{itemize}
\item 규칙에 묶인 구간은 계획에서 정한 대로 짝지은 다시 뽑기 400번에서 나와요. 그 몬테카를로 오차는 판정을 가른 거리에 견주면 작아요. 예외는 수정 F3이에요. 그 아래 끝은 허용 폭에 걸려 있고, 다른 시드 50개 가운데 하나에서는 아래 끝이 −0.01보다 위로 올라와 기준을 통과해요(부록~\ref{app:explore}). 규칙 B의 판정은 그 시드에 기대지 않아요. F2가 모든 시드에서 실패하기 때문이에요.
\item 판정은 세 조건의 교집합이라 다중 비교 보정이 필요 없어요. 하지만 판정 옆에 보고한 많은 대비는 명목상이라, 한 묶음의 시험으로 읽으면 안 돼요.
\item 계약을 다시 뽑는 부트스트랩은 스펜더들을 서로 독립으로 다뤄요. 그런데 새 주인 하나가 여러 스펜더를 허락할 수 있어요.
\item 트레이더 코호트는 작아요. 스펜더에서 가져온 허용 폭은 수정 분석에서는 트레이더가 겨우 닿는 경계에, 등록 분석에서는 닿지 않는 곳에 있었어요(\ref{sec:precision}절).
\end{itemize}

\subsection{외적 타당도}

증거는 한 롤업, 주인 셋 이상에게 허락을 받은 그 롤업의 계약, 동결 직전 1년으로 만든 점수(등록 분석에서는 2023년 10월부터의 모든 사건), 석 달짜리 라벨 기간 두 개에서 나왔어요. 각 그래프는 코호트의 한 걸음 에고 네트워크라서, 나가는 무게와 걷는 길은 체인 전체가 아니라 그 에고 네트워크의 것이에요. 보증이 한 걸음보다 멀리 퍼지는 더 깊은 표본에서는 다르게 움직일 수 있어요. 시험한 것은 균등 순간 이동을 쓰는 고정된 전체 무게뿐이에요. 점마다 다른 무게나 학습한 무게, 정상 벡터의 혼합, 층 사이에 점의 복사본을 두는 방식, 층마다 다른 감쇠 계수는 열린 문제로 남아 있어요.
